% -----------------------------------------------------------------------------
% Komponen yang dipakai lintas berkas modul.
% -----------------------------------------------------------------------------

% --- penanda tipografis ------------------------------------------------------
% \berkas   : nama berkas atau direktori
% \perintah : perintah shell atau nama objek basis data
% \istilah  : istilah asing yang sengaja dipertahankan
\newcommand{\berkas}[1]{\texttt{#1}}
\newcommand{\perintah}[1]{\texttt{#1}}
\newcommand{\istilah}[1]{\emph{#1}}

% --- logo dan spasi huruf pada sampul ---------------------------------------
\newcommand{\sampullogo}[2]{%
  \IfFileExists{../03_branding/#1}{%
    \includegraphics[height=#2]{../03_branding/#1}%
  }{\IfFileExists{assets/#1}{\includegraphics[height=#2]{assets/#1}}{}}%
}
% Spasi huruf .05em seperti pada berkas SVG sampul.
\newcommand{\sampulspasi}[1]{\textls[50]{\MakeUppercase{#1}}}

% -----------------------------------------------------------------------------
% Kotak berlabel.
% -----------------------------------------------------------------------------
\newtcolorbox{capaian}[1][]{
  enhanced, breakable, colback=DWOrange!5, colframe=DWOrange!55,
  boxrule=.7pt, arc=2pt, left=8pt, right=8pt, top=6pt, bottom=6pt,
  fonttitle=\bfseries\footnotesize\color{white},
  title={CAPAIAN PEMBELAJARAN MODUL}, colbacktitle=DWOrange, #1}

\newtcolorbox{konsep}[2][]{
  enhanced, breakable, colback=DWSage!8, colframe=DWSage!70,
  boxrule=.7pt, arc=2pt, left=8pt, right=8pt, top=6pt, bottom=6pt,
  fonttitle=\bfseries\footnotesize\color{white},
  title={KONSEP: \MakeUppercase{#2}}, colbacktitle=DWSage!85!DWInk, #1}

\newtcolorbox{luaran}[1][]{
  enhanced, breakable, colback=DWButter!35, colframe=DWPeach!85!DWInk,
  boxrule=.7pt, arc=2pt, left=8pt, right=8pt, top=6pt, bottom=6pt,
  fonttitle=\bfseries\footnotesize\color{DWInk},
  title={LUARAN YANG DIKUMPULKAN}, colbacktitle=DWPeach, coltitle=DWInk, #1}

\newtcolorbox{checkpoint}[1][]{
  enhanced, breakable, colback=SuccessGreen!6, colframe=SuccessGreen!60,
  boxrule=.7pt, arc=2pt, left=8pt, right=8pt, top=6pt, bottom=6pt,
  fonttitle=\bfseries\footnotesize\color{white},
  title={CHECKPOINT --- DIVERIFIKASI ASISTEN}, colbacktitle=SuccessGreen, #1}

\newtcolorbox{peringatan}[1][]{
  enhanced, breakable, colback=SignalRed!6, colframe=SignalRed!60,
  boxrule=.7pt, arc=2pt, left=8pt, right=8pt, top=6pt, bottom=6pt,
  fonttitle=\bfseries\footnotesize\color{white},
  title={KESALAHAN YANG LAZIM MUNCUL}, colbacktitle=SignalRed, #1}

\newtcolorbox{catatan}[1][]{
  enhanced, breakable, colback=SoftGray, colframe=LineGray!45,
  boxrule=.4pt, arc=2pt, left=8pt, right=8pt, top=6pt, bottom=6pt,
  fonttitle=\bfseries\footnotesize\color{Ink},
  title={CATATAN}, colbacktitle=LineGray!25, coltitle=Ink, #1}

\newtcolorbox{catatanasisten}[1][]{
  enhanced, breakable, colback=Ink!4, colframe=Ink!50,
  boxrule=.7pt, arc=2pt, left=8pt, right=8pt, top=6pt, bottom=6pt,
  fonttitle=\bfseries\footnotesize\color{white},
  title={CATATAN UNTUK ASISTEN}, colbacktitle=Ink, #1}

% -----------------------------------------------------------------------------
% Penomoran latihan, terhitung ulang pada setiap bab.
% -----------------------------------------------------------------------------
\newcounter{latihan}[chapter]
\newcommand{\latihan}[1]{%
  \refstepcounter{latihan}%
  \needspace{6\baselineskip}%
  \subsection*{%
    \colorbox{DWOrange}{\textcolor{white}{\ \strut Latihan \thechapter.\thelatihan\ }}%
    \enspace\color{DWOrange}#1}%
  \vspace{.2em}%
}

% -----------------------------------------------------------------------------
% Judul langkah praktikum, terlindung dari judul yatim di kaki halaman.
% Argumen kedua adalah perkiraan waktu pengerjaan.
% -----------------------------------------------------------------------------
\newcommand{\langkah}[2]{%
  \needspace{8\baselineskip}%
  \subsubsection*{#1}%
  \hfill{\footnotesize\color{Slate}\itshape[#2]}\par\vspace{-.4em}%
}

% -----------------------------------------------------------------------------
% Tabel identitas modul -- dipakai identik oleh sepuluh modul.
% Baris ditulis sebagai "\textbf{Komponen} & Keterangan \\".
% Baris CPMK mengacu pada pemetaan pada CPMK.md melalui kuliah penopang.
% -----------------------------------------------------------------------------
\newenvironment{identitasmodul}{%
  \begin{center}\small
  \begin{tabular}{@{}p{0.26\textwidth}>{\raggedright\arraybackslash}p{0.66\textwidth}@{}}
  \toprule
  \textbf{Komponen} & \textbf{Keterangan} \\
  \midrule
}{%
  \bottomrule
  \end{tabular}
  \end{center}
}

% -----------------------------------------------------------------------------
% Tabel alur 120 menit -- dipakai identik oleh sepuluh modul.
% -----------------------------------------------------------------------------
\newenvironment{alurwaktu}{%
  \begin{center}\small
  \begin{tabular}{@{}p{0.14\textwidth}p{0.78\textwidth}@{}}
  \toprule
  \textbf{Menit} & \textbf{Kegiatan} \\
  \midrule
}{%
  \bottomrule
  \end{tabular}
  \end{center}
}

% -----------------------------------------------------------------------------
% Rujukan ke slide kuliah dan ke sumber pustaka.
% -----------------------------------------------------------------------------
\newcommand{\slideref}[1]{%
  {\footnotesize\color{Slate}Materi kuliah: Pertemuan~#1.\par}}
\newcommand{\acuan}[1]{%
  \par\vspace{.25em}{\footnotesize\color{MidGray}Acuan: #1\par}}

% -----------------------------------------------------------------------------
% Checklist pengumpulan.
% -----------------------------------------------------------------------------
\newlist{ceklis}{itemize}{1}
\setlist[ceklis]{label=$\square$,leftmargin=1.4em,itemsep=.3em}
